\documentclass[11pt]{article}
\usepackage{fontspec}
\setmainfont{Latin Modern Roman}
\usepackage{amsmath,amssymb,mathtools,unicode-math}
\setmathfont{Latin Modern Math}
\usepackage{xurl,hyperref}
\hypersetup{hidelinks,pdftitle={Problem statement}}
\newcommand{\sref}[2]{\hyperlink{source:#1:#2}{[S#2]}}
\newcommand{\cataloguescope}[1]{\paragraph{Mathematical question.}#1}
\begin{document}
\section{Erdős–Faber–Lovász conjecture}\label{statement}

\subsection{Definitions and mathematical statement}

\paragraph{Shared notation.}$\mathbb N=\{1,2,\ldots\}$, and $\mathbb Z,\mathbb Q,\mathbb R,\mathbb C$ denote the integers, rationals, reals and complex numbers. Any other conventions are specified locally.


Let $k\ge1$, and let $G$ be the union of $k$ complete subgraphs $Q_1,\ldots,Q_k$, each on $k$ vertices, such that their edge sets are pairwise disjoint. Thus $V(G)=\bigcup_iV(Q_i)$ and $E(G)=\bigcup_iE(Q_i)$; different $Q_i$ may share vertices. Let $\chi(G)$ be the proper vertex chromatic number.
\paragraph{Statement recorded in the supplied source.}
Must every such finite simple graph satisfy $\chi(G)=k$? \sref{statement}{2}
\subsection{Short English statement}
Can the union of k edge-disjoint cliques of order k always be colored with k colors?
\subsection{Sources}
\begin{description}
\item[\hypertarget{source:statement:1}{S1}] ULAM.AI and the UnsolvedMath Contributors, UnsolvedMath: A Curated Collection of Open Mathematics Problems, catalogue entry. CC BY 4.0. \url{https://huggingface.co/datasets/ulamai/UnsolvedMath}.
\item[\hypertarget{source:statement:2}{S2}] Open Problem Garden, Erdős–Faber–Lovász conjecture, original node 55984, statement and mathematical discussion. \url{https://www.openproblemgarden.org/op/erdos_faber_lovasz_conjecture}.
\end{description}
\label{end:statement}
\end{document}
